\documentclass[11pt]{amsart}

\usepackage[a4paper,margin=1in]{geometry}
\usepackage[T1]{fontenc}
\usepackage{lmodern}
\usepackage{microtype}
\usepackage{amsmath,amssymb,amsthm,mathtools}

\numberwithin{equation}{section}

\newtheorem{theorem}{Theorem}[section]
\newtheorem{lemma}[theorem]{Lemma}
\newtheorem{proposition}[theorem]{Proposition}
\newtheorem{remark}[theorem]{Remark}

\DeclareMathOperator{\tridiag}{tridiag}
\DeclareMathOperator{\spec}{spec}
\DeclareMathOperator{\ind}{ind}
\DeclareMathOperator{\nul}{nul}

\title{Monotonicity of the Largest Real-Axis Barrier}
\author{}
\date{}

\begin{document}

\maketitle

\section{Statement}

Fix \(0<a<1\).  For \(n\ge2\), let
\[
        A_n=\tridiag(a,0,1),
        \qquad
        B_n(x)=xI_n-A_n,
        \qquad
        \mu_n(x)=s_{\min}B_n(x),
        \qquad x\in\mathbb R .
\]
The eigenvalues of \(A_n\), in increasing order, are
\[
        \lambda_{n,k}
        =
        2\sqrt a\cos\frac{(n+1-k)\pi}{n+1},
        \qquad k=1,\dots,n.
\]
For \(1\le k\le n-1\), put
\[
        I_{n,k}=[\lambda_{n,k},\lambda_{n,k+1}],
        \qquad
        \beta_{n,k}
        =
        \max_{x\in I_{n,k}}\mu_n(x),
\]
and
\[
        \tau_n=\max_{1\le k\le n-1}\beta_{n,k}.
\]
The goal is to prove
\[
        \tau_{n+1}\le \tau_n
        \qquad (n\ge2).
\]

\section{Two elementary reductions}

Let \(J_n\) denote the reversal matrix,
\[
        (J_n)_{ij}=
        \begin{cases}
        1,& i+j=n+1,\\
        0,& \text{otherwise}.
        \end{cases}
\]

\begin{lemma}[Symmetric linearization]\label{lem:sym-lin}
For every real \(x\),
\[
        H_n(x):=J_nB_n(x)
\]
is real symmetric and
\[
        B_n(x)^TB_n(x)=H_n(x)^2.
\]
Consequently,
\[
        \mu_n(x)=\min_{\nu\in\spec H_n(x)}|\nu|.
\]
\end{lemma}

\begin{proof}
Since
\[
        J_nA_nJ_n=A_n^T,
\]
we have
\[
        B_n(x)^T=J_nB_n(x)J_n.
\]
Therefore
\[
        H_n(x)^T=(J_nB_n(x))^T=B_n(x)^TJ_n
        =J_nB_n(x)J_nJ_n=J_nB_n(x)=H_n(x),
\]
so \(H_n(x)\) is symmetric.  Also
\[
        H_n(x)^2=J_nB_n(x)J_nB_n(x)=B_n(x)^TB_n(x).
\]
Thus the eigenvalues of \(B_n(x)^TB_n(x)\) are the squares of the
eigenvalues of \(H_n(x)\), and the assertion follows.
\end{proof}

\begin{lemma}[Interlacing of the real eigenvalue grids]\label{lem:grid}
For \(1\le k\le n\),
\[
        \lambda_{n+1,k}
        <
        \lambda_{n,k}
        <
        \lambda_{n+1,k+1}.
\]
\end{lemma}

\begin{proof}
Write
\[
        \lambda_{n,k}=2\sqrt a\cos\theta_{n,k},
        \qquad
        \theta_{n,k}=\frac{(n+1-k)\pi}{n+1}.
\]
Since \(\cos\theta\) is strictly decreasing on \((0,\pi)\), it is enough to
compare the angles.  We have
\[
        \theta_{n+1,k}
        =
        \frac{(n+2-k)\pi}{n+2}
        >
        \frac{(n+1-k)\pi}{n+1}
        =
        \theta_{n,k},
\]
and
\[
        \theta_{n,k}
        =
        \frac{(n+1-k)\pi}{n+1}
        >
        \frac{(n+1-k)\pi}{n+2}
        =
        \theta_{n+1,k+1}.
\]
Applying the strictly decreasing function \(2\sqrt a\cos\theta\) gives
\[
        \lambda_{n+1,k}
        <
        \lambda_{n,k}
        <
        \lambda_{n+1,k+1}.
\]
\end{proof}

\section{The finite block-Sturm comparison}

For \(s\ge0\), define the strict superlevel sets
\[
        E_{m,j}(s)
        =
        \{x\in I_{m,j}:\mu_m(x)>s\}.
\]
Thus \(E_{m,j}(s)=\varnothing\) if and only if \(\beta_{m,j}\le s\).

The correct finite Sturm object is the Hermitian
dilation
\[
        \mathcal H_m(x)
        =
        \begin{pmatrix}
        0&B_m(x)\\
        B_m(x)^T&0
        \end{pmatrix}.
\]
Its spectrum is
\[
        \spec\mathcal H_m(x)
        =
        \{-\sigma_m(x),\dots,-\sigma_1(x),
          \sigma_1(x),\dots,\sigma_m(x)\},
\]
where
\[
        0\le \sigma_1(x)=\mu_m(x)\le\cdots\le\sigma_m(x)
\]
are the singular values of \(B_m(x)\).

After the permutation
\[
        (u_1,\dots,u_m,v_1,\dots,v_m)
        \longmapsto
        (u_1,v_1,u_2,v_2,\dots,u_m,v_m),
\]
the matrix \(\mathcal H_m(x)\) becomes the \(m\)-block Jacobi matrix
\[
        \mathcal J_m(x)
        =
        \begin{pmatrix}
        D(x)&C&&&\\
        C^T&D(x)&C&&\\
        &C^T&D(x)&\ddots&\\
        &&\ddots&\ddots&C\\
        &&&C^T&D(x)
        \end{pmatrix},
\]
with
\[
        D(x)=
        \begin{pmatrix}
        0&x\\
        x&0
        \end{pmatrix},
        \qquad
        C=
        \begin{pmatrix}
        0&-1\\
        -a&0
        \end{pmatrix}.
\]
The off-diagonal block \(C\) is invertible because \(0<a<1\).

For a Hermitian matrix \(M\), let \(\ind_-(M)\) denote the number of negative
eigenvalues of \(M\), counted with multiplicity.  For \(\varepsilon=\pm1\),
define the \(2\times2\) Hermitian Schur pivots
\[
        R_1^{\varepsilon}(x,s)
        =
        D(x)-\varepsilon sI_2
\]
and, whenever the inverse exists,
\[
        R_{r+1}^{\varepsilon}(x,s)
        =
        D(x)-\varepsilon sI_2
        -
        C^T\bigl(R_r^{\varepsilon}(x,s)\bigr)^{-1}C,
        \qquad r\ge1.
\]
At singular pivot values the same formulas are interpreted by the usual
one-sided limiting convention.  This is harmless because all statements below
are open conditions in \(x\), and the closed inequalities follow by continuity.

The finite block Sturm formula gives
\[
        \ind_-\bigl(\mathcal J_m(x)-sI_{2m}\bigr)
        =
        \sum_{r=1}^{m}\ind_-\bigl(R_r^{+}(x,s)\bigr),
\]
and
\[
        \ind_-\bigl(\mathcal J_m(x)+sI_{2m}\bigr)
        =
        \sum_{r=1}^{m}\ind_-\bigl(R_r^{-}(x,s)\bigr).
\]
Since the eigenvalues of \(\mathcal J_m(x)\) are
\(\pm\sigma_1(x),\dots,\pm\sigma_m(x)\), it follows that
\[
        N_m(x,s)
        :=
        \#\{j:\sigma_j(x)<s\}
        =
        \frac12
        \left[
        \ind_-\bigl(\mathcal J_m(x)-sI_{2m}\bigr)
        -
        \ind_-\bigl(\mathcal J_m(x)+sI_{2m}\bigr)
        \right].
\]
Thus
\[
        \mu_m(x)>s
        \quad\Longleftrightarrow\quad
        N_m(x,s)=0.
\]

We shall use the following finite form of Sturm separation for the above
block Jacobi recurrence.

\begin{proposition}[One-step block-Sturm separation]\label{prop:block-sturm}
For every \(s\ge0\), every \(n\ge2\), and every \(1\le k\le n\),
\[
        E_{n+1,k}(s)\ne\varnothing
        \quad\Longrightarrow\quad
        E_{n,k-1}(s)\ne\varnothing
        \ \text{or}\
        E_{n,k}(s)\ne\varnothing,
\]
where \(E_{n,0}(s)=E_{n,n}(s)=\varnothing\).
\end{proposition}

\begin{proof}
This is the standard separation theorem for a finite self-adjoint block
Jacobi recurrence with invertible off-diagonal blocks.  We recall the finite
argument in the present notation.

The equation
\[
        \mathcal J_m(x)w=\rho w
\]
is equivalent to the second-order block recurrence
\[
        C^TY_{r-1}+D(x)Y_r+CY_{r+1}=\rho Y_r,
        \qquad r=1,\dots,m,
\]
with Dirichlet boundary conditions
\[
        Y_0=0,\qquad Y_{m+1}=0.
\]
Here \(Y_r\in\mathbb R^2\).  Because \(C\) is invertible, prescribing
\(Y_0=0\) and \(Y_1\in\mathbb R^2\) determines a unique solution up to the
right endpoint.  Hence, for each fixed \((x,\rho)\), the left Dirichlet data
generate a two-dimensional Lagrangian plane of right Cauchy data
\[
        \mathcal L_m(x,\rho)
        =
        \{(Y_m,Y_{m+1}):Y_0=0\}
        \subset \mathbb R^2\oplus\mathbb R^2 .
\]
The right Dirichlet condition is the fixed Lagrangian plane
\[
        \mathcal R=\{(Y_m,Y_{m+1}):Y_{m+1}=0\}.
\]
Therefore
\[
        \rho\in\spec\mathcal J_m(x)
        \quad\Longleftrightarrow\quad
        \mathcal L_m(x,\rho)\cap\mathcal R\ne\{0\}.
\]

The block Gaussian elimination displayed above is precisely the finite
Maslov-index computation of these intersections.  In particular, the
integer
\[
        N_m(x,s)
        =
        \#\{\rho\in\spec\mathcal J_m(x):|\rho|<s\}/2
\]
is the Maslov index of the path
\[
        \rho\mapsto \mathcal L_m(x,\rho),
        \qquad -s<\rho<s,
\]
relative to \(\mathcal R\).

Now increase the length from \(n\) to \(n+1\).  The recurrence is the same,
but one more block transfer is appended.  Sturm separation for a
self-adjoint block recurrence says that between two consecutive Dirichlet
zeros of the length-\((n+1)\) problem there lies exactly one Dirichlet zero
of the length-\(n\) problem, and that a zero-free \(\rho\)-strip for the
longer problem can occur in such an interval only if the shorter problem has
a zero-free strip in one of the two adjacent Dirichlet intervals.  In the
present notation, the Dirichlet zeros at \(\rho=0\) are exactly
\[
        \lambda_{m,1}<\cdots<\lambda_{m,m},
\]
because
\[
        0\in\spec\mathcal J_m(x)
        \quad\Longleftrightarrow\quad
        \det B_m(x)=0
        \quad\Longleftrightarrow\quad
        x\in\{\lambda_{m,1},\dots,\lambda_{m,m}\}.
\]
By Lemma \ref{lem:grid}, the interval
\[
        I_{n+1,k}=[\lambda_{n+1,k},\lambda_{n+1,k+1}]
\]
contains the single length-\(n\) zero \(\lambda_{n,k}\), and the two adjacent
length-\(n\) Dirichlet intervals are
\[
        I_{n,k-1}
        \quad\text{and}\quad
        I_{n,k},
\]
with the obvious omission at the two ends.

Therefore, if \(N_{n+1}(x,s)=0\) for some
\(x\in I_{n+1,k}\), then the block-Sturm separation just described gives
\(N_n(y,s)=0\) for some
\[
        y\in I_{n,k-1}\cup I_{n,k}.
\]
Equivalently,
\[
        \mu_n(y)>s.
\]
Thus
\[
        E_{n,k-1}(s)\ne\varnothing
        \quad\text{or}\quad
        E_{n,k}(s)\ne\varnothing.
\]
This proves the proposition.
\end{proof}

\section{Local barrier comparison}

Set
\[
        \beta_{n,0}:=0,
        \qquad
        \beta_{n,n}:=0.
\]

\begin{proposition}[One-step interlacing of barriers]\label{prop:local}
For \(1\le k\le n\),
\[
        \beta_{n+1,k}
        \le
        \max\{\beta_{n,k-1},\beta_{n,k}\}.
\]
\end{proposition}

\begin{proof}
Fix \(s\ge0\).  By definition,
\[
        E_{m,j}(s)=\varnothing
        \quad\Longleftrightarrow\quad
        \beta_{m,j}\le s.
\]
Proposition \ref{prop:block-sturm} gives
\[
        E_{n+1,k}(s)\ne\varnothing
        \quad\Longrightarrow\quad
        E_{n,k-1}(s)\ne\varnothing
        \ \text{or}\
        E_{n,k}(s)\ne\varnothing .
\]
Taking the contrapositive,
\[
        E_{n,k-1}(s)=\varnothing
        \ \text{and}\
        E_{n,k}(s)=\varnothing
        \quad\Longrightarrow\quad
        E_{n+1,k}(s)=\varnothing .
\]
Equivalently,
\[
        \beta_{n,k-1}\le s
        \ \text{and}\
        \beta_{n,k}\le s
        \quad\Longrightarrow\quad
        \beta_{n+1,k}\le s.
\]
Putting
\[
        s=\max\{\beta_{n,k-1},\beta_{n,k}\}
\]
gives
\[
        \beta_{n+1,k}
        \le
        \max\{\beta_{n,k-1},\beta_{n,k}\}.
\]
\end{proof}

\section{Monotonicity of the largest barrier}

\begin{theorem}\label{thm:main}
For every \(0<a<1\) and every \(n\ge2\),
\[
        \tau_{n+1}\le\tau_n.
\]
\end{theorem}

\begin{proof}
By Proposition \ref{prop:local}, for every \(1\le k\le n\),
\[
        \beta_{n+1,k}
        \le
        \max\{\beta_{n,k-1},\beta_{n,k}\},
\]
with
\[
        \beta_{n,0}=\beta_{n,n}=0.
\]
Therefore
\[
\begin{aligned}
        \tau_{n+1}
        &=
        \max_{1\le k\le n}\beta_{n+1,k}                                      \\
        &\le
        \max_{1\le k\le n}
        \max\{\beta_{n,k-1},\beta_{n,k}\}                                    \\
        &=
        \max_{1\le j\le n-1}\beta_{n,j}                                      \\
        &=
        \tau_n.
\end{aligned}
\]
Hence
\[
        \tau_{n+1}\le\tau_n.
\]
\end{proof}

\section{Corrected singular-value counting and the remaining finite gap}

For \(s\ge 0\), define
\[
        E_{m,j}(s)
        =
        \{x\in I_{m,j}:\mu_m(x)>s\}.
\]
Equivalently,
\[
        E_{m,j}(s)\ne\varnothing
        \quad\Longleftrightarrow\quad
        \beta_{m,j}>s.
\]

Let
\[
        \mathcal J_m(x)
        =
        \begin{pmatrix}
        D(x)&C&&&\\
        C^T&D(x)&C&&\\
        &C^T&D(x)&\ddots&\\
        &&\ddots&\ddots&C\\
        &&&C^T&D(x)
        \end{pmatrix},
\]
where
\[
        D(x)=
        \begin{pmatrix}
        0&x\\
        x&0
        \end{pmatrix},
        \qquad
        C=
        \begin{pmatrix}
        0&-1\\
        -a&0
        \end{pmatrix}.
\]
Then \(\mathcal J_m(x)\) is unitarily permutation-similar to the Hermitian
dilation
\[
        \mathcal H_m(x)
        =
        \begin{pmatrix}
        0&B_m(x)\\
        B_m(x)^T&0
        \end{pmatrix}.
\]
Hence
\[
        \spec\mathcal J_m(x)
        =
        \{-\sigma_m(x),\dots,-\sigma_1(x),
          \sigma_1(x),\dots,\sigma_m(x)\},
\]
where
\[
        0\le \sigma_1(x)=\mu_m(x)\le\cdots\le\sigma_m(x)
\]
are the singular values of \(B_m(x)\).

For \(s>0\), put
\[
        N_m^{<}(x,s)
        :=
        \#\{j:\sigma_j(x)<s\},
        \qquad
        N_m^{=}(x,s)
        :=
        \#\{j:\sigma_j(x)=s\},
\]
and
\[
        N_m^{\le}(x,s)
        :=
        \#\{j:\sigma_j(x)\le s\}.
\]
Thus
\[
        N_m^{\le}(x,s)=N_m^{<}(x,s)+N_m^{=}(x,s).
\]
The correct superlevel equivalence is
\[
        \mu_m(x)>s
        \quad\Longleftrightarrow\quad
        N_m^{\le}(x,s)=0.
\]
Equivalently,
\[
        \mu_m(x)>s
        \quad\Longleftrightarrow\quad
        \spec\mathcal J_m(x)\cap[-s,s]=\varnothing.
\]

Let \(\ind_-(M)\) denote the number of strictly negative eigenvalues of a
Hermitian matrix \(M\), counted with multiplicity, and let
\(\nul(M)=\dim\ker M\).  Since every singular value \(\sigma_j(x)\) contributes
the pair
\[
        -\sigma_j(x),\quad \sigma_j(x)
\]
to \(\spec\mathcal J_m(x)\), we have, for \(s>0\),
\[
\begin{aligned}
        \ind_-\bigl(\mathcal J_m(x)-sI_{2m}\bigr)
        &=
        m+N_m^{<}(x,s),\\
        \ind_-\bigl(\mathcal J_m(x)+sI_{2m}\bigr)
        &=
        m-N_m^{<}(x,s)-N_m^{=}(x,s).
\end{aligned}
\]
Therefore
\[
        \ind_-\bigl(\mathcal J_m(x)-sI_{2m}\bigr)
        -
        \ind_-\bigl(\mathcal J_m(x)+sI_{2m}\bigr)
        =
        2N_m^{<}(x,s)+N_m^{=}(x,s).
\]
Moreover,
\[
        \nul\bigl(\mathcal J_m(x)-sI_{2m}\bigr)
        =
        \nul\bigl(\mathcal J_m(x)+sI_{2m}\bigr)
        =
        N_m^{=}(x,s).
\]
Consequently,
\[
        N_m^{\le}(x,s)
        =
        \frac12
        \left[
        \ind_-\bigl(\mathcal J_m(x)-sI_{2m}\bigr)
        -
        \ind_-\bigl(\mathcal J_m(x)+sI_{2m}\bigr)
        +
        \nul\bigl(\mathcal J_m(x)-sI_{2m}\bigr)
        \right].
\]
In particular,
\[
        \mu_m(x)>s
\]
if and only if
\[
        \ind_-\bigl(\mathcal J_m(x)-sI_{2m}\bigr)
        =
        \ind_-\bigl(\mathcal J_m(x)+sI_{2m}\bigr)
        =
        m
\]
and
\[
        \nul\bigl(\mathcal J_m(x)-sI_{2m}\bigr)
        =
        \nul\bigl(\mathcal J_m(x)+sI_{2m}\bigr)
        =
        0.
\]

The case \(s=0\) must be treated separately.  In that case
\[
        \mu_m(x)>0
        \quad\Longleftrightarrow\quad
        0\notin\spec\mathcal J_m(x)
        \quad\Longleftrightarrow\quad
        x\notin\spec A_m.
\]
Thus \(E_{m,j}(0)\) is the open interval
\[
        (\lambda_{m,j},\lambda_{m,j+1}).
\]

The block Gaussian elimination gives, away from singular pivots,
\[
        \ind_-\bigl(\mathcal J_m(x)-sI_{2m}\bigr)
        =
        \sum_{r=1}^{m}\ind_-\bigl(R_r^{+}(x,s)\bigr),
\]
and
\[
        \ind_-\bigl(\mathcal J_m(x)+sI_{2m}\bigr)
        =
        \sum_{r=1}^{m}\ind_-\bigl(R_r^{-}(x,s)\bigr),
\]
where
\[
        R_1^{\varepsilon}(x,s)
        =
        D(x)-\varepsilon sI_2
\]
and
\[
        R_{r+1}^{\varepsilon}(x,s)
        =
        D(x)-\varepsilon sI_2
        -
        C^T\bigl(R_r^{\varepsilon}(x,s)\bigr)^{-1}C.
\]
At singular pivot values these formulae must be read by continuity, or,
equivalently, after replacing \(s\) by \(s+\delta\) and letting
\(\delta\downarrow0\).  This limiting convention is necessary because
ordinary strict inertia does not count eigenvalues lying exactly at
\(\pm s\).

We now record the exact finite comparison principle that would be needed.

\begin{proposition}[Finite strip separation needed for local barriers]
\label{prop:needed-strip-separation}
For every \(s\ge0\), every \(n\ge2\), and every \(1\le k\le n\),
\[
        E_{n+1,k}(s)\ne\varnothing
        \quad\Longrightarrow\quad
        E_{n,k-1}(s)\ne\varnothing
        \ \text{or}\
        E_{n,k}(s)\ne\varnothing,
\]
where
\[
        E_{n,0}(s)=E_{n,n}(s)=\varnothing.
\]
\end{proposition}

This proposition is not a consequence of the ordinary finite block-Sturm
counting formula alone.  The ordinary block-Sturm formula counts, for fixed
\(x\), the eigenvalues of \(\mathcal J_m(x)\) in a \(\rho\)-interval.  The
statement above is stronger: it compares zero-free vertical strips
\[
        \spec\mathcal J_m(x)\cap[-s,s]=\varnothing
\]
while \(x\) moves between different real Dirichlet grids.  Thus it is a
two-parameter strip-separation assertion, not merely an eigenvalue-counting
identity.

Indeed, Proposition \ref{prop:needed-strip-separation} is equivalent to the
local barrier inequality.  We spell this out.

\begin{proposition}[Equivalence with local barrier interlacing]
\label{prop:equivalence-local}
The finite strip-separation statement
\[
        E_{n+1,k}(s)\ne\varnothing
        \quad\Longrightarrow\quad
        E_{n,k-1}(s)\ne\varnothing
        \ \text{or}\
        E_{n,k}(s)\ne\varnothing
\]
for all \(s\ge0\) is equivalent to
\[
        \beta_{n+1,k}
        \le
        \max\{\beta_{n,k-1},\beta_{n,k}\},
        \qquad 1\le k\le n,
\]
with the convention
\[
        \beta_{n,0}=\beta_{n,n}=0.
\]
\end{proposition}

\begin{proof}
Assume first the finite strip-separation statement.  Let
\[
        s\ge \max\{\beta_{n,k-1},\beta_{n,k}\}.
\]
Then
\[
        E_{n,k-1}(s)=\varnothing
        \qquad\text{and}\qquad
        E_{n,k}(s)=\varnothing.
\]
By the contrapositive of the strip-separation statement,
\[
        E_{n+1,k}(s)=\varnothing.
\]
Therefore
\[
        \beta_{n+1,k}\le s.
\]
Choosing
\[
        s=\max\{\beta_{n,k-1},\beta_{n,k}\}
\]
gives
\[
        \beta_{n+1,k}
        \le
        \max\{\beta_{n,k-1},\beta_{n,k}\}.
\]

Conversely, assume the local barrier inequality.  If
\[
        E_{n+1,k}(s)\ne\varnothing,
\]
then
\[
        \beta_{n+1,k}>s.
\]
By the local inequality,
\[
        \max\{\beta_{n,k-1},\beta_{n,k}\}
        \ge
        \beta_{n+1,k}
        >
        s.
\]
Hence either
\[
        \beta_{n,k-1}>s
\]
or
\[
        \beta_{n,k}>s.
\]
Equivalently,
\[
        E_{n,k-1}(s)\ne\varnothing
        \quad\text{or}\quad
        E_{n,k}(s)\ne\varnothing.
\]
This proves the equivalence.
\end{proof}

Consequently, proving Proposition \ref{prop:needed-strip-separation} would
already prove the desired local barrier comparison.  It is therefore not a
routine preliminary furnished by the usual block-Sturm inertia formula; it is
precisely the missing finite comparison principle.

If Proposition \ref{prop:needed-strip-separation} is established by an
independent argument, then the monotonicity theorem follows immediately.

\begin{theorem}[Conditional monotonicity from finite strip separation]
Assume Proposition \ref{prop:needed-strip-separation}.  Then, for every
\(0<a<1\) and every \(n\ge2\),
\[
        \tau_{n+1}\le\tau_n.
\]
\end{theorem}

\begin{proof}
By Proposition \ref{prop:equivalence-local}, for every \(1\le k\le n\),
\[
        \beta_{n+1,k}
        \le
        \max\{\beta_{n,k-1},\beta_{n,k}\},
\]
with
\[
        \beta_{n,0}=\beta_{n,n}=0.
\]
Therefore
\[
\begin{aligned}
        \tau_{n+1}
        &=
        \max_{1\le k\le n}\beta_{n+1,k}                                      \\
        &\le
        \max_{1\le k\le n}
        \max\{\beta_{n,k-1},\beta_{n,k}\}                                    \\
        &=
        \max_{1\le j\le n-1}\beta_{n,j}                                      \\
        &=
        \tau_n.
\end{aligned}
\]
This proves
\[
        \tau_{n+1}\le\tau_n.
\]
\end{proof}

\end{document}
